\documentclass[11pt,a4paper]{article}
\usepackage[a4paper,margin=1.8cm]{geometry}
\usepackage[T1]{fontenc}
\usepackage[utf8]{inputenc}
\usepackage{lmodern}
\usepackage{microtype}
\usepackage{amsmath,amssymb}
\usepackage{mathastext}
\usepackage{upgreek}
\usepackage{enumitem}
\usepackage{tabularx}
\usepackage{array}
\usepackage{xcolor}
\usepackage{needspace}
\usepackage{ragged2e}

\definecolor{jawab}{HTML}{166534}
\definecolor{jelas}{HTML}{374151}

\pagestyle{empty}
\setlength{\parindent}{0pt}
\setlength{\parskip}{4pt}
\setlist[enumerate]{leftmargin=*,itemsep=2pt,topsep=3pt}
\renewcommand{\arraystretch}{1.25}
\setlength{\arrayrulewidth}{0.8pt}
\setlength{\tabcolsep}{8pt}

\newcommand{\answer}[1]{\par\medskip\textbf{\color{jawab}ANSWER:} #1\par}
\newcommand{\explanation}[1]{\textbf{\color{jelas}EXPLANATION:} #1\par\medskip}
\newcommand{\question}[1]{\Needspace{8\baselineskip}\par\medskip\textbf{#1}\quad}
\newenvironment{choices}{\begin{enumerate}[label=\Alph*.,leftmargin=1.8em]}{\end{enumerate}}

\begin{document}

\question{1.} Himpunan penyelesaian persamaan
\[
|3x+2|=11
\]
adalah ....
\begin{choices}
\item $\left\{-\frac{13}{3},3\right\}$
\item $\left\{-3,\frac{13}{3}\right\}$
\item $\left\{-\frac{11}{3},3\right\}$
\item $\left\{-\frac{13}{3},\frac{11}{3}\right\}$
\item $\{-13,3\}$
\end{choices}
\answer{A}
\explanation{Untuk $|f(x)|=c$ dengan $c\ge0$, berlaku $f(x)=c$ atau $f(x)=-c$. Maka $3x+2=11$ memberi $x=3$, sedangkan $3x+2=-11$ memberi $x=-\frac{13}{3}$.}

\question{2.} Pada garis bilangan, titik $x$ berjarak 7 satuan dari titik 4. Nilai $x$ yang mungkin adalah ....
\begin{choices}
\item $-11$ dan 3
\item $-7$ dan 11
\item $-3$ dan 11
\item 3 dan 7
\item 4 dan 11
\end{choices}
\answer{C}
\explanation{Pernyataan ``$x$ berjarak 7 satuan dari 4'' ditulis $|x-4|=7$. Jadi $x-4=7$ atau $x-4=-7$, sehingga $x=11$ atau $x=-3$.}

\question{3.} Himpunan penyelesaian dari
\[
|2x-5|=|x+4|
\]
adalah ....
\begin{choices}
\item $\left\{-9,\frac13\right\}$
\item $\left\{\frac13,9\right\}$
\item $\{-1,9\}$
\item $\left\{-\frac13,9\right\}$
\item $\left\{-9,-\frac13\right\}$
\end{choices}
\answer{B}
\explanation{Karena $|A|=|B|$, maka $A=B$ atau $A=-B$. Dari $2x-5=x+4$ diperoleh $x=9$. Dari $2x-5=-(x+4)$ diperoleh $3x=1$, sehingga $x=\frac13$.}

\question{4.} Himpunan penyelesaian persamaan
\[
|2x+1|=x+7
\]
adalah ....
\begin{choices}
\item $\left\{-\frac83,6\right\}$
\item $\left\{-6,\frac83\right\}$
\item $\{-8,6\}$
\item $\left\{-\frac83\right\}$
\item $\{6\}$
\end{choices}
\answer{A}
\explanation{Ruas kanan harus non-negatif: $x+7\ge0$, sehingga $x\ge-7$. Cabang pertama $2x+1=x+7$ memberi $x=6$. Cabang kedua $2x+1=-(x+7)$ memberi $3x=-8$, sehingga $x=-\frac83$. Keduanya memenuhi syarat $x\ge-7$.}

\question{5.} Banyaknya penyelesaian real persamaan
\[
|x^2-5x|=6
\]
adalah ....
\begin{choices}
\item 0
\item 1
\item 2
\item 3
\item 4
\end{choices}
\answer{E}
\explanation{Persamaan setara dengan $x^2-5x=6$ atau $x^2-5x=-6$. Persamaan pertama memberi $(x-6)(x+1)=0$, yaitu $x=6,-1$. Persamaan kedua memberi $(x-2)(x-3)=0$, yaitu $x=2,3$. Ada 4 penyelesaian real berbeda.}

\question{6.} Persamaan
\[
|x-2|=k
\]
memiliki tepat satu penyelesaian real. Nilai $k$ adalah ....
\begin{choices}
\item $-2$
\item $-1$
\item 0
\item 1
\item 2
\end{choices}
\answer{C}
\explanation{Jika $k<0$, tidak ada solusi. Jika $k>0$, ada dua solusi $x=2-k$ dan $x=2+k$. Tepat satu solusi terjadi saat $k=0$, yaitu $x=2$.}

\question{7.} Untuk setiap $x,y\in\mathbb{R}$ dengan $y\ne0$, pernyataan yang benar adalah ....
\begin{choices}
\item $|x|\ge0$
\item $|x|=|-x|$
\item $|xy|=|x|+|y|$
\item $\left|\frac{x}{y}\right|=\frac{|x|}{|y|}$
\item $|x|^2=x^2$
\end{choices}
\answer{A, B, D, E}
\explanation{Sifat nilai mutlak yang berlaku adalah non-negatif, simetri, sifat pembagian, dan $|x|^2=x^2$. Untuk perkalian berlaku $|xy|=|x|\,|y|$, bukan penjumlahan.}

\question{8.} Diketahui persamaan
\[
|x+1|=3.
\]
Pernyataan yang benar adalah ....
\begin{choices}
\item Salah satu penyelesaiannya adalah $x=-4$.
\item Salah satu penyelesaiannya adalah $x=2$.
\item Jumlah kedua penyelesaian adalah $-2$.
\item Hasil kali kedua penyelesaian adalah 8.
\item Kedua penyelesaian masing-masing berjarak 3 satuan dari $-1$.
\end{choices}
\answer{A, B, C, E}
\explanation{Dari $|x+1|=3$ diperoleh $x+1=3$ atau $x+1=-3$, sehingga $x=2$ atau $x=-4$. Jumlahnya $-2$, hasil kalinya $-8$, dan keduanya berjarak 3 dari pusat $-1$.}

\question{9.} Diberikan persamaan
\[
|x-1|=2x+5.
\]
Pernyataan yang benar adalah ....
\begin{choices}
\item Syarat ruas kanan adalah $x\ge-\frac52$.
\item $x=-6$ merupakan penyelesaian persamaan.
\item $x=-\frac43$ merupakan penyelesaian persamaan.
\item Himpunan penyelesaiannya hanya memiliki satu anggota.
\item Persamaan tidak memiliki penyelesaian real.
\end{choices}
\answer{A, C, D}
\explanation{Syarat $2x+5\ge0$ memberi $x\ge-\frac52$. Cabang $x-1=2x+5$ menghasilkan $x=-6$, tetapi gagal syarat. Cabang $x-1=-(2x+5)$ menghasilkan $3x=-4$, yaitu $x=-\frac43$, dan nilai ini memenuhi syarat.}

\question{10.} Diketahui $\upalpha<\upbeta$ merupakan dua penyelesaian persamaan
\[
|x-5|=2.
\]
Bandingkan kuantitas P dan Q berikut.
\begin{center}
\begin{tabular}{|>{\centering\arraybackslash}p{0.35\textwidth}|>{\centering\arraybackslash}p{0.35\textwidth}|}
\hline
\textbf{P} & \textbf{Q} \\
\hline
$\upbeta-\upalpha$ & 4 \\
\hline
\end{tabular}
\end{center}
\begin{choices}
\item $P>Q$
\item $P<Q$
\item $P=Q$
\item Hubungan $P$ dan $Q$ tidak dapat ditentukan.
\end{choices}
\answer{C}
\explanation{Penyelesaian $|x-5|=2$ adalah $x=3$ dan $x=7$. Karena $\upalpha=3$ dan $\upbeta=7$, maka $P=7-3=4=Q$.}

\question{11.} Diketahui $\upalpha$ dan $\upbeta$ merupakan penyelesaian persamaan
\[
|2x-3|=5.
\]
Bandingkan P dan Q berikut.
\begin{center}
\begin{tabular}{|>{\centering\arraybackslash}p{0.35\textwidth}|>{\centering\arraybackslash}p{0.35\textwidth}|}
\hline
\textbf{P} & \textbf{Q} \\
\hline
$\upalpha\upbeta$ & $-3$ \\
\hline
\end{tabular}
\end{center}
\begin{choices}
\item $P>Q$
\item $P<Q$
\item $P=Q$
\item Hubungan $P$ dan $Q$ tidak dapat ditentukan.
\end{choices}
\answer{B}
\explanation{Dari $2x-3=5$ diperoleh $x=4$, sedangkan dari $2x-3=-5$ diperoleh $x=-1$. Jadi $P=(4)(-1)=-4$, sedangkan $Q=-3$. Maka $P<Q$.}

\question{12.} Bilangan real $x$ memenuhi
\[
|x+2|=4.
\]
Bandingkan P dan Q berikut.
\begin{center}
\begin{tabular}{|>{\centering\arraybackslash}p{0.35\textwidth}|>{\centering\arraybackslash}p{0.35\textwidth}|}
\hline
\textbf{P} & \textbf{Q} \\
\hline
$x^2$ & 9 \\
\hline
\end{tabular}
\end{center}
\begin{choices}
\item $P>Q$
\item $P<Q$
\item $P=Q$
\item Hubungan $P$ dan $Q$ tidak dapat ditentukan.
\end{choices}
\answer{D}
\explanation{Persamaan menghasilkan $x=2$ atau $x=-6$. Jika $x=2$, maka $P=4<Q$. Jika $x=-6$, maka $P=36>Q$. Hubungan tidak dapat ditentukan secara tunggal.}

\question{13.} Diketahui persamaan
\[
|x-a|=5.
\]
Berapakah nilai $a$?

(1) Salah satu penyelesaiannya adalah $x=8$.\\
(2) Salah satu penyelesaian lainnya adalah $x=-2$.
\begin{choices}
\item Pernyataan (1) saja cukup, tetapi (2) saja tidak cukup.
\item Pernyataan (2) saja cukup, tetapi (1) saja tidak cukup.
\item Kedua pernyataan bersama-sama cukup, tetapi masing-masing tidak cukup.
\item Masing-masing pernyataan cukup.
\item Kedua pernyataan tidak cukup.
\end{choices}
\answer{C}
\explanation{Dari (1), $|8-a|=5$ memberi $a=3$ atau $a=13$, sehingga belum cukup. Dari (2), $|-2-a|=5$ memberi $a=3$ atau $a=-7$, juga belum cukup. Jika keduanya benar, nilai yang konsisten hanya $a=3$.}

\question{14.} Diketahui persamaan
\[
|x-b|=0.
\]
Berapakah himpunan penyelesaiannya?

(1) $b=4$.\\
(2) Persamaan mempunyai tepat satu penyelesaian real.
\begin{choices}
\item Pernyataan (1) saja cukup, tetapi (2) saja tidak cukup.
\item Pernyataan (2) saja cukup, tetapi (1) saja tidak cukup.
\item Kedua pernyataan bersama-sama cukup, tetapi masing-masing tidak cukup.
\item Masing-masing pernyataan cukup.
\item Kedua pernyataan tidak cukup.
\end{choices}
\answer{A}
\explanation{Persamaan $|x-b|=0$ selalu memiliki satu solusi, yaitu $x=b$. Pernyataan (1) memberi $b=4$, sehingga HP $=\{4\}$. Pernyataan (2) tidak memberi nilai $b$, jadi tidak cukup menentukan HP secara numerik.}

\question{15.} Apakah persamaan
\[
|x+1|=k
\]
mempunyai dua penyelesaian real yang berbeda?

(1) $k>0$.\\
(2) $k\ge0$ dan $k\ne0$.
\begin{choices}
\item Pernyataan (1) saja cukup, tetapi (2) saja tidak cukup.
\item Pernyataan (2) saja cukup, tetapi (1) saja tidak cukup.
\item Kedua pernyataan bersama-sama cukup, tetapi masing-masing tidak cukup.
\item Masing-masing pernyataan cukup.
\item Kedua pernyataan tidak cukup.
\end{choices}
\answer{D}
\explanation{Persamaan $|x+1|=k$ memiliki dua solusi berbeda tepat ketika $k>0$. Pernyataan (1) langsung menjamin hal itu. Pernyataan (2) juga setara dengan $k>0$. Jadi masing-masing cukup.}

\end{document}
